
\section{Procedural RF Simulation with HOUSIONNA}\label{sec:simulator}
\subsection{Scene Representation and Export}\label{sec:scene}Houdini generates geometry, assigns material attributes, and places gateways, sensor nodes, and measurement surfaces. A parameterised Houdini node network can represent alternative crop layouts or infrastructure configurations without manually rebuilding every scene. The exporter groups triangular geometry by object name and writes meshes, a scene description, material parameters, device definitions, and solver settings.

Export checks include triangulation, object names, material consistency, and device definitions. These are checks of simulation input integrity; they do not establish biological realism or material accuracy. A preview stage loads the exported scene in Sionna and returns its geometry for inspection in Houdini. This helps identify coordinate, scale, and export discrepancies before interpreting propagation results.

The current implementation uses Sionna RT 2.2.0 and a single isotropic, vertically polarised antenna at each gateway and sensor node~\cite{housionna}. The solver accepts controls for line of sight, specular and diffuse reflection, transmission/refraction, diffraction, interaction depth, and sampling budgets. Enabling an interaction does not guarantee that a chosen scene representation models the corresponding canopy physics accurately.

\gilles{What can be chaned/variable? + pics?}



\begin{figure*}[tp]
\centering
\newcommand{\gridimg}[2]{%
  \begin{minipage}[t]{0.19\linewidth}\raggedright
    \includegraphics[width=\linewidth]{figures/procedural/#1_#2}\\
    {\scriptsize x=#1\\y=#2\par}
  \end{minipage}}%
\newcommand{\gridrow}[1]{%
  \gridimg{0}{#1}\hfill\gridimg{1}{#1}\hfill\gridimg{2}{#1}\hfill\gridimg{3}{#1}\hfill\gridimg{4}{#1}\par\smallskip}%
\gridrow{0}
\gridrow{1}
\gridrow{2}
\caption{Procedural scene variations.}
\label{fig:procedural_matrix}
\end{figure*}

\begin{figure*}[t]
\centering
\newcommand{\icon}[1]{\includegraphics[height=8.5mm]{figures/icons/#1.pdf}}
\definecolor{houdini}{HTML}{E07B39}\definecolor{sionna}{HTML}{2F9E44}\definecolor{browser}{HTML}{2563EB}\definecolor{analysis}{HTML}{7C3AED}\definecolor{outputline}{HTML}{6B7280}
\begin{tikzpicture}[
  font=\scriptsize,
  box/.style={draw=#1,fill=#1!7,rounded corners=3pt,align=center,text width=2.3cm,
              minimum height=1.95cm,inner sep=2pt,line width=0.9pt,
              execute at begin node={\hyphenpenalty=10000\exhyphenpenalty=10000}},
  box/.default=gray,
  todo/.style={box=gray,dashed,fill=white},
  output/.style={box=outputline,fill=white,rounded corners=0pt,line width=0.7pt},
  arr/.style={-{Stealth[length=1.8mm,width=1.4mm]},line width=0.7pt,black!80},
  todoarr/.style={arr,dashed},
  dot/.style={circle,fill=black!80,inner sep=0pt,minimum size=1.3mm},
  stage/.style={fill=black!5,rounded corners=5pt},
  stagelabel/.style={font=\scriptsize\bfseries,anchor=north west,inner sep=4pt,text=black!70,align=left},
  swatch/.style={draw=#1,fill=#1!7,line width=0.9pt,rounded corners=1.5pt,minimum width=3.2mm,minimum height=2.4mm},
  % The whole box links to the section that describes it.
  link/.style={alias=linked,append after command={
    let \p1=(linked.north west),\p2=(linked.south east),\n1={\x2-\x1},\n2={\y1-\y2} in
    node[inner sep=0pt,outer sep=0pt,anchor=north west] at (\p1)
      {\hyperref[#1]{\phantom{\rule{\n1}{\n2}}}}}},
]
% Stages: what happens (the box colour shows which tool does it)
\fill[stage] (-1.35,2.0) rectangle (1.35,-3.7);
\fill[stage] (1.6,2.0)   rectangle (7.3,-1.25);
\fill[stage] (7.55,2.0)  rectangle (13.2,-3.7);
\node[stagelabel] at (-1.35,2.0) {1 Procedural\\scene generation};
\node[stagelabel] at (1.6,2.0)   {2 Ray tracing};
\node[stagelabel] at (7.55,2.0)  {3 Evaluation};

% 1 Procedural scene generation
\node[box=houdini,link=sec:materials] (proc) at (0,0)
  {\icon{scene}\\[1pt]\textbf{Procedural scene}\\crops, greenhouse, devices};
\node[todo,link=sec:twin] (cal) at (0,-2.55)
  {\icon{calibrate}\\[1pt]\textbf{Calibrate}\\materials and antennas};
% 2 Ray tracing
\node[box=houdini,link=sec:scene] (exp) at (2.95,0)
  {\icon{export}\\[1pt]\textbf{Check and export}\\geometry and materials};
\node[box=sionna,link=sec:execution] (rt) at (5.9,0)
  {\icon{raytrace}\\[1pt]\textbf{Ray tracing}\\on a GPU, in batches};
% 3 Evaluation
\node[output,link=sec:returned] (out) at (8.85,0)
  {\icon{results}\\[1pt]\textbf{Results}\\paths, links, radio maps};
\node[box=browser,link=sec:returned] (view) at (11.8,0)
  {\icon{viewer}\\[1pt]\textbf{Web viewer}\\browse and publish runs};
\node[box=analysis,link=sec:metrics] (eng) at (8.85,-2.55)
  {\icon{engine}\\[1pt]\textbf{Application engine}\\coverage, RF power;\\Python or other tools};
\node[box=houdini,link=sec:returned] (hou) at (11.8,-2.55)
  {\icon{houdini}\\[1pt]\textbf{Inspect in Houdini}\\results on the scene};
% Towards a digital twin, between the stages
\node[todo,link=sec:twin] (meas) at (4.42,-2.55)
  {\icon{measure}\\[1pt]\textbf{Field measurements}\\RF and crop state};

\draw[arr] (proc)--(exp);
\draw[arr] (exp)--(rt);
\draw[arr] (rt)--(out);
\draw[arr] (out)--(view);
\draw[arr] (out)--(eng);
\draw[arr] (8.85,-1.42) node[dot]{} -| (hou.north);
% Iteration back to the scene, around the left of the stages
\draw[arr] (view.north) -- (11.8,2.5) -- node[above,inner sep=2pt]
  {\includegraphics[height=2.4mm]{figures/icons/iterate.pdf}\enspace iterate: change the scene or sweep parameters}
  (-1.65,2.5) -- (-1.65,0) -- (proc.west);
% Field measurements are compared with the predictions and calibrate the scene
\draw[todoarr] (meas) -- node[above,inner sep=2pt]{compare} (eng);
\draw[todoarr] (meas) -- node[above,inner sep=2pt]{calibrate} (cal);
\draw[todoarr] (cal) -- node[left,inner sep=2pt]{update} (proc);

% Legend: the tool behind each box
\node[anchor=north west,inner sep=0pt,align=left] at (13.5,2.0) {\textbf{Tool}\\[3pt]
  \tikz\node[swatch=houdini]{};\enspace Houdini\\[3pt]
  \tikz\node[swatch=sionna]{};\enspace Sionna RT\\[3pt]
  \tikz\node[swatch=browser]{};\enspace Browser\\[3pt]
  \tikz\node[swatch=analysis]{};\enspace Python or other\\[6pt]
  \tikz\node[draw=outputline,fill=white,line width=0.7pt,minimum width=3.2mm,minimum height=2.4mm]{};\enspace output data\\[3pt]
  \tikz\node[swatch=gray,dashed,fill=white]{};\enspace not yet implemented};
\end{tikzpicture}
\caption{The \housionna{} workflow in three stages; the colour of each box shows the tool that performs the step, and the square white box is output data. Dashed elements are needed for an agricultural digital twin, in which field measurements calibrate the materials and antennas of the scene. Each box links to the section that describes it.}
\label{fig:pipeline}
\end{figure*}




\subsection{Execution and Reproducibility}\label{sec:execution}The solver runs in a separate Python environment, locally on an NVIDIA GPU or a CPU under Linux, or on a remote GPU host. Results are written as compressed NumPy arrays and imported into Houdini as attributed geometry. This separation allows the modelling workstation and compute machine to differ.

For parameter sweeps, Houdini task-operator (TOP) work items, one per parameter combination, are exported to individual run folders. A background batch process solves work items sequentially within one Sionna session. Successful items can be skipped when their recorded solver settings match, and failed items are recorded for later retry.

Run records include the exported scene, resolved material properties, devices, settings, a solver copy, and software provenance. Batch records also preserve work-item parameters. Recording these assets makes the simulated geometry inspectable even after the procedural node network has changed. The records do not yet include the random seed of the path solver or hardware details, which a complete release should add.

\subsection{Returned Quantities}\label{sec:returned}Propagation paths are imported as polylines carrying gateway and sensor node identifiers, complex amplitudes, delays, interaction types, and object associations. Per-link outputs include path counts, coherent and incoherent gains, RMS delay spread, and frequency-response magnitudes. Radio maps provide path gain per gateway over triangular measurement surfaces. A web viewer built on three.js displays the same paths, radio maps, and link metrics without a Houdini installation, steps through the runs of a batch, and can be published as a read-only website.


The distinction between link and map outputs matters. Coherent link quantities depend on relative path phase and can vary over small spatial displacements. Radio maps average power over cells and do not retain the phase information required for coherent \gls{wpt} from multiple gateways. Nor does a map of the strongest gateway represent the combined field of several synchronised gateways. Per-gateway maps can nevertheless be combined: with maximum-ratio transmission, the RF power at the sensor node equals the total gateway RF power times the sum of the per-gateway gains, so its local average follows from the maps without phase information. \Cref{sec:metrics} uses this to evaluate strategies for multiple gateways over an area, and the link outputs to evaluate them at individual sensor nodes.

\subsection{Agricultural Geometry and Radio Materials}\label{sec:materials}
The scene parameters include plant spacing, canopy height and depth, leaf orientations, infrastructure layout, and device positions; gateways and sensor nodes can be placed anywhere, including on plants and on mobile robots. Plants are generated with L-systems in three growth stages, optionally with fruit, so growth is represented by a sequence of static scenes. People, autonomous robots, windows, and metal frames can be added as disturbances of the channel. A geometric growth parameter is not automatically a validated physiological model.



\begin{figure}[t]
\centering
\newcommand{\stageimg}[1]{%
  \begin{minipage}[t]{0.32\linewidth}\raggedright
    \includegraphics[width=\linewidth]{figures/growth_stage/#1}\\[-1pt]
    {\scriptsize #1\par}
  \end{minipage}}%
\stageimg{seedling}\hfill\stageimg{vegetative}\hfill\stageimg{mature}
\caption{Plant growth stages.}
\label{fig:growth_stages}
\end{figure}



The built-in ITU-R P.2040 materials of Sionna are bounded in frequency~\cite{itu2040}: apart from wood and glass, they are defined only from \SI{1}{GHz}, which excludes the \SI{868}{MHz} and \SI{915}{MHz} bands. \housionna{} therefore adds frequency-dependent models and records a source and status for each material: Cole--Cole fits of common building materials measured from \SIrange{0.2}{67}{GHz}~\cite{zhekov2020}, conductors for metal structures, water and nutrient solution~\cite{liebe1991}, soils~\cite{itu527}, and the dual-dispersion vegetation model for leaves and stems~\cite{ulaby1987}. The vegetation coefficients are not yet verified, and fruit and growing media use placeholder values, so plant materials must be treated as provisional. ITU-R P.833 provides an independent vegetation-attenuation reference for appropriate scenarios~\cite{itu833}; it does not replace calibration for the specific crop and geometry.

Sionna models interactions with surfaces: each surface a ray crosses acts as a slab with the full thickness of its radio material. A closed canopy shell, internal sheets, and explicit leaf geometry can produce different depth-dependent attenuation. Refining the mesh alone does not resolve this modelling choice. When a leaf or stem is small relative to wavelength, a surface ray approximation may miss relevant scattering behaviour. Sionna RT also discards paths containing a segment shorter than \SI{1}{cm}, so every path through a closed leaf or stem thinner than this is lost. Both effects were confirmed with analytical reference cases (free space, a single slab, a closed block, and thin glass) included in the repository~\cite{housionna}. Representation and material uncertainty must be assessed together.



